% Explanation certificate for heart-s1-c1 (genome a2df22ce-983d-47b6-bbd6-3176d38ed5ff).
% Generated by explaneat export (schema 1.4.0) from the JSON record alone;
% content sha256 555774df08fc19cc0dd5fe846e6c3e4fb3823536efda28def15a994d5545f6e2.
% Drafter recorded in the JSON: ai_agent (claude). Do not edit by hand: re-render with
%   uv run python -m explaneat export render --from heart-s1-c1.json
% TikZ modules are input from \explaneatcertdir{<label>} (empty by default: same
% directory; index.tex sets it to the genome's subdirectory).
\providecommand{\explaneatcertdir}[1]{}
\begin{drafted}[appendix: heart-s1-c1 explanation certificate]
\section{Explanation certificate: heart-s1-c1}\label{app:heart-s1-c1:sec}
This section is the explanation certificate of the heart-s1-c1 network. It is rendered from the canonical JSON record \texttt{heart-s1-c1.json} (schema 1.4.0, SHA-256 \texttt{555774df08fc19cc}\ldots), which also holds every evidence payload, the full graphs with weights, and the explanation record. \Cref{app:heart-s1-c1:summary} summarises the network and its coverage.

\begin{table}[htbp]
\centering\small
\begin{tabularx}{\linewidth}{@{}lX@{}}
\toprule
Dataset & Heart Disease (binary, prepared) \\
Experiment & 10seed-heart-s1, seed 1 \\
Genome & \texttt{a2df22ce-983d-47b6-bbd6-3176d38ed5ff} (NEAT key 9674, generation 39) \\
Explanation & \texttt{492866e6-a73b-4465-92a8-ecad82a2adcc} (35 operations) \\
Train AUC / accuracy & 0.9392 / 0.8802 (n = 242) \\
Test AUC / accuracy & 0.861 / 0.7705 (n = 61) \\
Output activation & sigmoid \\
Phenotype & 34 nodes (26 inputs, 7 hidden, 1 output), 30 enabled edges \\
Explained model state & 38 nodes, 30 enabled edges \\
Refactoring operations & split\_\allowbreak{}node 3 \\
Function preservation & max $|\Delta\,\mathrm{output}|$ = $3.33\times 10^{-16}$ on the training split \\
Annotations & 6 (5 leaf, 1 composition) \\
Structural coverage & 1 (37 of 37 candidate nodes; 9 Null) \\
Compositional coverage & 1 \\
Uncovered nodes / edges & 0 / 0 \\
Missing compositions & 0 \\
Full coverage & yes \\
Drafted by & AI-assisted draft \\
Verified by & unverified \\
Record hash (SHA-256) & \texttt{555774df08fc19cc}\ldots \\
\bottomrule
\end{tabularx}
\caption{Explanation certificate summary: heart-s1-c1.}
\label{app:heart-s1-c1:summary}
\end{table}

\paragraph{Drafter's summary.} Essentially logistic regression (all hidden units linearised: train AUC 0.9355 vs 0.9392; L2 logistic regression 0.9355). 13 direct edges carry logit SD 4.25 of 5.63; ca enters by three parallel routes (effective +2.60 logit per SD), cp:3 only via always-on relay 1173. The one consequential nonlinearity is gate 1801 = ReLU(1.54 restecg:0 + 0.82 oldpeak - 1.52): ST depression adds risk only with a normal resting ECG (+0.0034 AUC over linear). Typical angina (cp:1) lowers risk via switch 3210, less with more vessels; a small hinge lets higher max heart rate lower risk. Same design as the 10-seed heart-s1 record (gen-31 genome), with an added typical-angina switch and heart-rate hinge.

\paragraph{What resisted explanation.} Nothing resisted coverage; only input fan-out needed splitting (ca x3, slope:2, thalach). Four of seven hidden ReLUs (664, 1173, 2245, 2650) never clip on train, so their annotations describe affine relays (reweighting of ca, slope:2, cp:3 and the heart-rate hinge) rather than distinct concepts. 3210 is a pure cp:1 switch (on for all 22 typical-angina patients, off for the rest), so its best linear fit is almost exact. Nine unconnected inputs (sex:1, fbs:0, restecg:1, restecg:2, exang:0, exang:1, slope:1, thal:4, thal:7) listed as zero-effect root entries.

\subsection{Plain-language summary}\label{app:heart-s1-c1:plain}
\noindent{\small\textit{AI-assisted draft for review (unverified). Each sentence cites the claims (C01, C02, \ldots) that back it; every cited claim was computed on the model and the training split and passed its automatic checks. Magnitudes are changes in predicted risk for a typical patient; confidence words follow the test-split check and n. The text describes what the model does, not the world.}}\par\smallskip
\paragraph{Summary (drafted).} For a typical patient (a 56-year-old man with no chest pain symptoms and median test results) the model predicts a risk of heart disease of about 18\%; it adds up most of its inputs independently, and five parts add to that sum: two add more for each extra vessel coloured, one lowers the risk for typical angina, one lowers it for non-anginal chest pain, one acts only on the 42 of 242 training patients with a normal resting ECG and marked ST depression, and one lowers it for a higher maximum heart rate, with nothing until about 145 beats per minute, then steadily.\,{\scriptsize\textcolor{NavyBlue}{[C29, C30, C40, C31, C43, C38, C46]}} These add up independently, and these lower the risk: being a woman takes a typical patient's predicted risk from about 18\% to below 1\% and fasting blood sugar above 120 mg/dL to about 8\%; through this sum alone, compared with no chest pain symptoms, typical angina or non-anginal chest pain lowers it to about 2\% and atypical angina to about 5\% (chest pain type also acts through the parts below).\,{\scriptsize\textcolor{NavyBlue}{[C10, C19, C12, C17]}} A reversible defect (or the undefined code 4) on the thallium scan raises it to about 55\%; higher cholesterol (from about 8\% at 193 mg/dL to about 45\% at 308 mg/dL), higher resting blood pressure (from about 9\% at 110 mm Hg to about 35\% at 155 mm Hg) and older age (from about 10\% at 42 years to about 25\% at 66 years) raise it steadily.\,{\scriptsize\textcolor{NavyBlue}{[C05, C07, C08, C13]}} Through this sum alone, each extra vessel coloured by fluoroscopy adds about the same (from about 18\% at 0 to about 80\% at 2 vessels) and a flat ST segment raises it to about 87\%, while the maximum heart rate barely changes it; all three also act through the parts below.\,{\scriptsize\textcolor{NavyBlue}{[C04, C03, C21]}} Rarer results also count (the model has little data on this, n = 15 and 17): a fixed defect on the thallium scan lowers the risk to about 3\%, and through this sum alone a downsloping ST segment raises it to about 32\%; the model ignores exercise-induced angina.\,{\scriptsize\textcolor{NavyBlue}{[C14, C15, C28]}} This part combines chest pain type with the number of vessels coloured: through it, typical angina lowers a typical patient's predicted risk from about 18\% to about 2\%, and each extra vessel adds about the same, from about 18\% at 0 to about 61\% at 2 vessels.\,{\scriptsize\textcolor{NavyBlue}{[C31, C30]}} Across the whole model, more vessels raise the predicted risk, and even more so with typical angina, each extra vessel adding about the same: going from 0 to 2 vessels moves a typical patient from below 1\% to about 55\% with typical angina, and from about 18\% to about 99\% with any other chest pain type.\,{\scriptsize\textcolor{NavyBlue}{[C32]}} Across the whole model, typical angina takes a typical patient's predicted risk to below 1\%.\,{\scriptsize\textcolor{NavyBlue}{[C09]}} This part combines the resting ECG with ST depression on exercise and acts on 42 of 242 training patients, nearly all with a normal resting ECG, 26 of whom had heart disease; for the typical patient it acts on, the predicted risk is about 98\% without it and about 99\% with it.\,{\scriptsize\textcolor{NavyBlue}{[C38]}} Through it, more ST depression raises the predicted risk only once it rises above about 1.0 mm, for patients with a normal resting ECG, in an uneven way: for a typical such patient from about 13\% at 0 mm to about 63\% at 3.5 mm.\,{\scriptsize\textcolor{NavyBlue}{[C37]}} More ST depression therefore raises the predicted risk, but only if the resting ECG is normal: going from 0 to 2.8 mm moves a typical patient from about 18\% to about 55\% with a normal resting ECG, in an uneven way, and leaves it at about 18\% otherwise.\,{\scriptsize\textcolor{NavyBlue}{[C36]}} This part combines the number of vessels coloured with a flat ST segment: through it, each extra vessel adds about the same, from about 18\% at 0 to about 36\% at 2 vessels for a typical patient, and a flat ST segment raises the risk from about 18\% to about 23\%.\,{\scriptsize\textcolor{NavyBlue}{[C40, C41]}} Across the whole model, each extra vessel takes a typical patient from about 18\% at 0 to about 99\% at 2 vessels, about the same for each, flattening out near the top because the predicted risk cannot go much higher, and a flat ST segment takes it to about 91\%.\,{\scriptsize\textcolor{NavyBlue}{[C01, C02]}} This part reads only non-anginal chest pain: through it, non-anginal chest pain lowers a typical patient's predicted risk from about 18\% to about 8\%.\,{\scriptsize\textcolor{NavyBlue}{[C43]}} Across the whole model, non-anginal chest pain takes it to about 1\%.\,{\scriptsize\textcolor{NavyBlue}{[C11]}} This part reads only the maximum heart rate: a higher maximum heart rate lowers the predicted risk with nothing until about 145 beats per minute, then steadily, for a typical patient from about 19\% at 118 to about 14\% at 175 beats per minute.\,{\scriptsize\textcolor{NavyBlue}{[C45, C46]}} Across the whole model the same holds, from about 19\% at 118 to about 14\% at 175 beats per minute.\,{\scriptsize\textcolor{NavyBlue}{[C20]}}
\paragraph{The whole model.} The model adds up its inputs into one score: the vessels coloured (each extra vessel adding about the same), the ST segment, chest pain type and sex each move the predicted risk on their own, five parts add to that sum, and exercise-induced angina is ignored.\,{\scriptsize\textcolor{NavyBlue}{[C29, C01, C02, C10, C28]}} \textit{Evidence.} It is made up of the inputs that feed the score directly and five parts: vessels with a flat ST segment, the non-anginal chest pain relay, ST depression with a normal ECG, the maximum heart rate part and typical angina with vessels (each explained below); the score is their sum, which was checked on all 242 training patients.\,{\scriptsize\textcolor{NavyBlue}{[C29]}} Changing one input at a time across the whole model, each extra vessel adds about the same, taking a typical patient's predicted risk from about 18\% at 0 to about 99\% at 2 vessels, flattening out near the top because the predicted risk cannot go much higher, against about 80\% through the sum alone.\,{\scriptsize\textcolor{NavyBlue}{[C01, C04]}} More ST depression raises it in an uneven way across the whole model, from about 18\% at 0 mm to about 55\% at 2.8 mm; it enters the score only through the part with a normal ECG.\,{\scriptsize\textcolor{NavyBlue}{[C06]}} The model ignores exercise-induced angina.\,{\scriptsize\textcolor{NavyBlue}{[C28]}}
\paragraph{Typical angina with vessels.} This part lowers the predicted risk for typical angina and adds risk for each extra vessel coloured, about the same for each, with the vessels counting more for patients with typical angina.\,{\scriptsize\textcolor{NavyBlue}{[C31, C30, C32]}} \textit{Evidence.} Through it, typical angina takes a typical patient's predicted risk from about 18\% to about 2\%, and each extra vessel adds about the same, from about 18\% at 0 to about 61\% at 2 vessels.\,{\scriptsize\textcolor{NavyBlue}{[C31, C30]}} Across the whole model, going from 0 to 2 vessels moves a typical patient from below 1\% to about 55\% with typical angina, each extra vessel adding about the same, and from about 18\% to about 99\% with any other chest pain type.\,{\scriptsize\textcolor{NavyBlue}{[C32]}} Its formula reproduces the part exactly on all 242 training patients.\,{\scriptsize\textcolor{NavyBlue}{[C33]}}
\paragraph{ST depression with a normal ECG.} This part raises the predicted risk with ST depression on exercise, in an uneven way, only for patients with a normal resting ECG, once it rises above about 1.0 mm.\,{\scriptsize\textcolor{NavyBlue}{[C37, C36]}} \textit{Evidence.} It acts on 42 of 242 training patients, nearly all with a normal resting ECG, 26 of whom had heart disease; for the typical patient it acts on, the predicted risk is about 98\% without it and about 99\% with it.\,{\scriptsize\textcolor{NavyBlue}{[C38]}} Through it, more ST depression raises the predicted risk only once it rises above about 1.0 mm, in an uneven way: for a typical patient with a normal resting ECG from about 13\% at 0 mm to about 63\% at 3.5 mm.\,{\scriptsize\textcolor{NavyBlue}{[C37]}} Going from 0 to 2.8 mm of ST depression moves a typical patient from about 18\% to about 55\% with a normal resting ECG, in an uneven way, and leaves it at about 18\% otherwise.\,{\scriptsize\textcolor{NavyBlue}{[C36]}} Its formula reproduces the part exactly on all 242 training patients.\,{\scriptsize\textcolor{NavyBlue}{[C39]}}
\paragraph{Vessels with flat ST segment.} This part adds risk for each extra vessel coloured, about the same for each, and for a flat ST segment.\,{\scriptsize\textcolor{NavyBlue}{[C40, C41]}} \textit{Evidence.} Through it, each extra vessel adds about the same, from about 18\% at 0 to about 36\% at 2 vessels for a typical patient, and a flat ST segment raises the risk from about 18\% to about 23\%.\,{\scriptsize\textcolor{NavyBlue}{[C40, C41]}} Its formula reproduces the part exactly on all 242 training patients.\,{\scriptsize\textcolor{NavyBlue}{[C42]}}
\paragraph{Non-anginal chest pain relay.} This part lowers the predicted risk for patients with non-anginal chest pain.\,{\scriptsize\textcolor{NavyBlue}{[C43]}} \textit{Evidence.} Through it, non-anginal chest pain takes a typical patient's predicted risk from about 18\% to about 8\%.\,{\scriptsize\textcolor{NavyBlue}{[C43]}} Its formula reproduces the part exactly on all 242 training patients.\,{\scriptsize\textcolor{NavyBlue}{[C44]}}
\paragraph{Maximum heart rate.} This part lowers the predicted risk for a higher maximum heart rate: nothing until about 145 beats per minute, then steadily.\,{\scriptsize\textcolor{NavyBlue}{[C45, C46]}} \textit{Evidence.} Through it, a higher maximum heart rate takes a typical patient from about 19\% at 71 to about 13\% at 182 beats per minute, with nothing until about 145 beats per minute, then steadily.\,{\scriptsize\textcolor{NavyBlue}{[C46]}} Its formula reproduces the part exactly on all 242 training patients.\,{\scriptsize\textcolor{NavyBlue}{[C47]}}
\paragraph{Worth checking against what is known.} Worth checking: patients with typical angina or non-anginal chest pain get the lowest predicted risk of all chest pain types (below 1\% and about 1\% for a typical patient), below patients with no chest pain symptoms.\,{\scriptsize\textcolor{NavyBlue}{[C09, C11]}} Worth checking: the model treats a fixed defect on the thallium scan as lowering the risk compared with a normal scan, from only 15 training patients.\,{\scriptsize\textcolor{NavyBlue}{[C14]}} Worth checking: the model treats fasting blood sugar above 120 mg/dL as lowering the risk of heart disease (for a typical patient from about 18\% to about 8\%).\,{\scriptsize\textcolor{NavyBlue}{[C19]}} Worth checking: going from 0 to 2 vessels coloured moves a typical patient's predicted risk from about 18\% to about 99\%.\,{\scriptsize\textcolor{NavyBlue}{[C01]}} Worth checking: ST depression on exercise raises the predicted risk only when the resting ECG is normal.\,{\scriptsize\textcolor{NavyBlue}{[C36]}} Worth checking: the model ignores exercise-induced angina.\,{\scriptsize\textcolor{NavyBlue}{[C28]}}
\paragraph{Claims.}
\par{\tiny\noindent
\begin{tabularx}{\linewidth}{@{}l>{\raggedright\arraybackslash}X>{\raggedright\arraybackslash}p{0.11\linewidth}>{\raggedright\arraybackslash}p{0.16\linewidth}>{\raggedright\arraybackslash}p{0.05\linewidth}>{\raggedright\arraybackslash}p{0.15\linewidth}l>{\raggedright\arraybackslash}p{0.08\linewidth}@{}}
\toprule
Claim & What is checked & Risk, typical patient & Estimate (log-odds; OR); range across patients & n & Test check & Verified & Words \\
\midrule
C01 & direction: number of vessels coloured by fluoroscopy: 0 vessels to 2 vessels (whole model) \newline\textit{shape: steady (equal steps); 100\% of patients same; verified; risk saturates near 100\%} & 18\% to 99\% & log-odds 5.89, OR 362; range across patients 5.83 to 5.83 & 242 & train; test unclear (AUC gain 0.041 [-0.032, 0.123], n=61) & \textcolor{OliveGreen}{yes} & on the training data only \\
C02 & direction: slope of the ST segment at peak exercise: an upsloping ST segment to a flat ST segment (whole model) \newline\textit{shape: categorical; highest to lowest: a flat ST segment, a downsloping ST segment, an upsloping ST segment} & 18\% to 91\% & log-odds 3.81, OR 44.9; range across patients 3.81 to 3.81 & 115, 110 & train; test disagrees (AUC gain -0.011 [-0.065, 0.036], n=61) & \textcolor{OliveGreen}{yes} & on the training data only \\
C03 & direction: slope of the ST segment at peak exercise: an upsloping ST segment to a flat ST segment (direct contribution only, not through any part) \newline\textit{shape: categorical; highest to lowest: a flat ST segment, a downsloping ST segment, an upsloping ST segment} & 18\% to 87\% & log-odds 3.46, OR 31.9; range across patients 3.46 to 3.46 & 115, 110 & train; test disagrees (AUC gain -0.012 [-0.059, 0.030], n=61) & \textcolor{OliveGreen}{yes} & on the training data only \\
C04 & direction: number of vessels coloured by fluoroscopy: 0 vessels to 2 vessels (direct contribution only, not through any part) \newline\textit{shape: steady (equal steps); 100\% of patients same; verified} & 18\% to 80\% & log-odds 2.91, OR 18.3; range across patients 2.91 to 2.91 & 242 & train; test disagrees (AUC gain -0.001 [-0.039, 0.032], n=61) & \textcolor{OliveGreen}{yes} & on the training data only \\
C05 & direction: thallium stress test result: a normal thallium scan to thallium stress test result other than a normal thallium scan, a fixed defect on the thallium scan (whole model) \newline\textit{shape: categorical; highest to lowest: thallium stress test result other than a normal thallium scan, a fixed defect on the thallium scan, a normal thallium scan, a fixed defect on the thallium scan} & 18\% to 55\% & log-odds 1.74, OR 5.71; range across patients 1.74 to 1.74 & 135, 92 & train; test unclear (AUC gain 0.012 [-0.017, 0.044], n=61) & \textcolor{OliveGreen}{yes} & on the training data only \\
C06 & direction: ST depression on exercise: 0.0 mm to 2.8 mm (whole model) \newline\textit{shape: hinge at 1.0 mm; 100\% of patients same; fallback; matches the fold at 1.0 mm} & 18\% to 55\% & log-odds 0.87, OR 2.38; range across patients 0.00 to 1.74 & 242 & train; test unclear (AUC gain 0.005 [-0.005, 0.022], n=61) & \textcolor{OliveGreen}{yes} & on the training data only \\
C07 & direction: serum cholesterol: 193 mg/dL (5.0 mmol/L) to 308 mg/dL (8.0 mmol/L) (whole model) \newline\textit{shape: steady; 100\% of patients same; verified} & 8\% to 45\% & log-odds 2.22, OR 9.23; range across patients 2.22 to 2.22 & 242 & train; test disagrees (AUC gain -0.022 [-0.066, 0.011], n=61) & \textcolor{OliveGreen}{yes} & on the training data only \\
C08 & direction: resting blood pressure: 110 mm Hg to 155 mm Hg (whole model) \newline\textit{shape: steady; 100\% of patients same; verified} & 9\% to 35\% & log-odds 1.65, OR 5.20; range across patients 1.65 to 1.65 & 242 & train; test unclear (AUC gain 0.000 [-0.022, 0.020], n=61) & \textcolor{OliveGreen}{yes} & on the training data only \\
C09 & direction: chest pain type: no chest pain (asymptomatic) to typical angina (whole model) \newline\textit{shape: categorical; highest to lowest: no chest pain (asymptomatic), atypical angina, non-anginal chest pain, typical angina} & 18\% to $<$1\% & log-odds -4.55, OR 0.01; range across patients -4.78 to -4.08 & 108, 22 & train; test agrees (AUC gain 0.041 [0.003, 0.083], n=61) & \textcolor{OliveGreen}{yes} & clearly \\
C10 & direction: sex: male to female (whole model) \newline\textit{shape: categorical; highest to lowest: male, female} & 18\% to $<$1\% & log-odds -3.94, OR 0.02; range across patients -3.94 to -3.94 & 159, 83 & train; test unclear (AUC gain 0.023 [-0.014, 0.068], n=61) & \textcolor{OliveGreen}{yes} & on the training data only \\
C11 & direction: chest pain type: no chest pain (asymptomatic) to non-anginal chest pain (whole model) \newline\textit{shape: categorical; highest to lowest: no chest pain (asymptomatic), atypical angina, non-anginal chest pain, typical angina} & 18\% to 1\% & log-odds -3.43, OR 0.03; range across patients -3.43 to -3.43 & 108, 70 & train; test agrees (AUC gain 0.041 [0.003, 0.083], n=61) & \textcolor{OliveGreen}{yes} & clearly \\
C12 & direction: chest pain type: no chest pain (asymptomatic) to chest pain type other than atypical angina, no chest pain (asymptomatic) (direct contribution only, not through any part) \newline\textit{shape: categorical; highest to lowest: no chest pain (asymptomatic), atypical angina, chest pain type other than atypical angina, no chest pain (asymptomatic)} & 18\% to 2\% & log-odds -2.51, OR 0.08; range across patients -2.51 to -2.51 & 108, 92 & train; test agrees (AUC gain 0.037 [0.006, 0.074], n=61) & \textcolor{OliveGreen}{yes} & clearly \\
\bottomrule
\end{tabularx}}
\par{\tiny\noindent
\begin{tabularx}{\linewidth}{@{}l>{\raggedright\arraybackslash}X>{\raggedright\arraybackslash}p{0.11\linewidth}>{\raggedright\arraybackslash}p{0.16\linewidth}>{\raggedright\arraybackslash}p{0.05\linewidth}>{\raggedright\arraybackslash}p{0.15\linewidth}l>{\raggedright\arraybackslash}p{0.08\linewidth}@{}}
\toprule
Claim & What is checked & Risk, typical patient & Estimate (log-odds; OR); range across patients & n & Test check & Verified & Words \\
\midrule
C13 & direction: age: 42 years to 66 years (whole model) \newline\textit{shape: steady; 100\% of patients same; verified} & 10\% to 25\% & log-odds 1.05, OR 2.87; range across patients 1.05 to 1.05 & 242 & train; test disagrees (AUC gain -0.008 [-0.018, 0.001], n=61) & \textcolor{OliveGreen}{yes} & on the training data only \\
C14 & direction: thallium stress test result: a normal thallium scan to a fixed defect on the thallium scan (whole model) \newline\textit{shape: categorical; highest to lowest: thallium stress test result other than a normal thallium scan, a fixed defect on the thallium scan, a normal thallium scan, a fixed defect on the thallium scan} & 18\% to 3\% & log-odds -1.83, OR 0.16; range across patients -1.83 to -1.83 & 135, 15 (small n) & train; test unclear (AUC gain 0.012 [-0.017, 0.044], n=61) & \textcolor{OliveGreen}{yes} & the model has little data on this; on the training data only \\
C15 & direction: slope of the ST segment at peak exercise: an upsloping ST segment to a downsloping ST segment (direct contribution only, not through any part) \newline\textit{shape: categorical; highest to lowest: a flat ST segment, a downsloping ST segment, an upsloping ST segment} & 18\% to 32\% & log-odds 0.77, OR 2.16; range across patients 0.77 to 0.77 & 115, 17 (small n) & train; test disagrees (AUC gain -0.012 [-0.059, 0.030], n=61) & \textcolor{OliveGreen}{yes} & the model has little data on this; on the training data only \\
C16 & direction: slope of the ST segment at peak exercise: an upsloping ST segment to a downsloping ST segment (whole model) \newline\textit{shape: categorical; highest to lowest: a flat ST segment, a downsloping ST segment, an upsloping ST segment} & 18\% to 32\% & log-odds 0.77, OR 2.16; range across patients 0.77 to 0.77 & 115, 17 (small n) & train; test disagrees (AUC gain -0.011 [-0.065, 0.036], n=61) & \textcolor{OliveGreen}{yes} & the model has little data on this; on the training data only \\
C17 & direction: chest pain type: no chest pain (asymptomatic) to atypical angina (direct contribution only, not through any part) \newline\textit{shape: categorical; highest to lowest: no chest pain (asymptomatic), atypical angina, chest pain type other than atypical angina, no chest pain (asymptomatic)} & 18\% to 5\% & log-odds -1.34, OR 0.26; range across patients -1.34 to -1.34 & 108, 42 & train; test agrees (AUC gain 0.037 [0.006, 0.074], n=61) & \textcolor{OliveGreen}{yes} & clearly \\
C18 & direction: chest pain type: no chest pain (asymptomatic) to atypical angina (whole model) \newline\textit{shape: categorical; highest to lowest: no chest pain (asymptomatic), atypical angina, non-anginal chest pain, typical angina} & 18\% to 5\% & log-odds -1.34, OR 0.26; range across patients -1.34 to -1.34 & 108, 42 & train; test agrees (AUC gain 0.041 [0.003, 0.083], n=61) & \textcolor{OliveGreen}{yes} & clearly \\
C19 & direction: fasting blood sugar: fasting blood sugar of 120 mg/dL or less to fasting blood sugar above 120 mg/dL (whole model) \newline\textit{shape: categorical; highest to lowest: fasting blood sugar of 120 mg/dL or less, fasting blood sugar above 120 mg/dL} & 18\% to 8\% & log-odds -0.92, OR 0.40; range across patients -0.92 to -0.92 & 209, 33 & train; test disagrees (AUC gain -0.001 [-0.010, 0.009], n=61) & \textcolor{OliveGreen}{yes} & on the training data only \\
C20 & direction: maximum heart rate: 118 beats per minute to 175 beats per minute (whole model) \newline\textit{shape: hinge at 145 beats per minute; 100\% of patients same; verified; matches the fold at 145 beats per minute} & 19\% to 14\% & log-odds -0.36, OR 0.70; range across patients -0.36 to -0.36 & 242 & train; test unclear (AUC gain 0.002 [0.000, 0.009], n=61) & \textcolor{OliveGreen}{yes} & slightly; on the training data only \\
C21 & direction: maximum heart rate: 118 beats per minute to 175 beats per minute (direct contribution only, not through any part) \newline\textit{shape: flat; 100\% of patients same; verified} & 18\% to 18\% & log-odds 0.005, OR 1.00; range across patients 0.005 to 0.005 & 242 & train; test unclear (AUC gain 0.000 [0.000, 0.000], n=61) & \textcolor{OliveGreen}{yes} & slightly; on the training data only \\
C22 & direction: resting ECG: a normal resting ECG to resting ECG other than a normal resting ECG (whole model) \newline\textit{shape: categorical; highest to lowest: a normal resting ECG, resting ECG other than a normal resting ECG} & 18\% to 18\% & log-odds -0.46, OR 0.63; range across patients -1.74 to 0.00 & 121, 121 & train; test disagrees (AUC gain -0.008 [-0.021, 0.000], n=61) & \textcolor{OliveGreen}{yes} & slightly; on the training data only \\
C23 & magnitude: number of vessels coloured by fluoroscopy: 0 vessels to 2 vessels, typical patient \newline\textit{shape: steady (equal steps); 100\% of patients same; verified; risk saturates near 100\%} & 18\% to 99\% & log-odds 5.83, OR 340 & 242 & train; test unclear (AUC gain 0.041 [-0.032, 0.123], n=61) & \textcolor{OliveGreen}{yes} & on the training data only \\
C24 & magnitude: slope of the ST segment at peak exercise: an upsloping ST segment to a flat ST segment, typical patient \newline\textit{shape: categorical; highest to lowest: a flat ST segment, a downsloping ST segment, an upsloping ST segment} & 18\% to 91\% & log-odds 3.81, OR 44.9 & 115, 110 & train; test disagrees (AUC gain -0.011 [-0.065, 0.036], n=61) & \textcolor{OliveGreen}{yes} & on the training data only \\
\bottomrule
\end{tabularx}}
\par{\tiny\noindent
\begin{tabularx}{\linewidth}{@{}l>{\raggedright\arraybackslash}X>{\raggedright\arraybackslash}p{0.11\linewidth}>{\raggedright\arraybackslash}p{0.16\linewidth}>{\raggedright\arraybackslash}p{0.05\linewidth}>{\raggedright\arraybackslash}p{0.15\linewidth}l>{\raggedright\arraybackslash}p{0.08\linewidth}@{}}
\toprule
Claim & What is checked & Risk, typical patient & Estimate (log-odds; OR); range across patients & n & Test check & Verified & Words \\
\midrule
C25 & magnitude: thallium stress test result: a normal thallium scan to thallium stress test result other than a normal thallium scan, a fixed defect on the thallium scan, typical patient \newline\textit{shape: categorical; highest to lowest: thallium stress test result other than a normal thallium scan, a fixed defect on the thallium scan, a normal thallium scan, a fixed defect on the thallium scan} & 18\% to 55\% & log-odds 1.74, OR 5.71 & 135, 92 & train; test unclear (AUC gain 0.012 [-0.017, 0.044], n=61) & \textcolor{OliveGreen}{yes} & on the training data only \\
C26 & magnitude: ST depression on exercise: 0.0 mm to 2.8 mm, typical patient \newline\textit{shape: hinge at 1.0 mm; 100\% of patients same; fallback; matches the fold at 1.0 mm} & 18\% to 55\% & log-odds 1.74, OR 5.68 & 242 & train; test unclear (AUC gain 0.005 [-0.005, 0.022], n=61) & \textcolor{OliveGreen}{yes} & on the training data only \\
C27 & magnitude: serum cholesterol: 193 mg/dL (5.0 mmol/L) to 308 mg/dL (8.0 mmol/L), typical patient \newline\textit{shape: steady; 100\% of patients same; verified} & 8\% to 45\% & log-odds 2.22, OR 9.23 & 242 & train; test disagrees (AUC gain -0.022 [-0.066, 0.011], n=61) & \textcolor{OliveGreen}{yes} & on the training data only \\
C28 & null: model ignores exercise-induced angina (exact rule) &  & max change 0.000 (max $|$change in log-odds$|$) & 242 & train & \textcolor{OliveGreen}{yes} &  \\
C29 & structure: made of ca\_\allowbreak{}slope2\_\allowbreak{}relay, cp3\_\allowbreak{}relay, st\_\allowbreak{}depression\_\allowbreak{}normal\_\allowbreak{}ecg\_\allowbreak{}gate,\allowbreak{} max\_\allowbreak{}heart\_\allowbreak{}rate\_\allowbreak{}hinge,\allowbreak{} typical\_\allowbreak{}angina\_\allowbreak{}ca\_\allowbreak{}switch and inputs age, sex, cp, trestbps, chol, fbs, thalach, slope, ca, thal; added &  & log-odds n/a, OR n/a & 242 & train & \textcolor{OliveGreen}{yes} &  \\
C30 & direction: number of vessels coloured by fluoroscopy: 0 vessels to 2 vessels (through this part) \newline\textit{shape: steady (equal steps); 100\% of patients same; verified} & 18\% to 61\% & log-odds 2.04, OR 7.73; range across patients 1.98 to 1.98 & 242 & train; test disagrees (AUC gain -0.003 [-0.032, 0.019], n=61) & \textcolor{OliveGreen}{yes} & on the training data only \\
C31 & direction: chest pain type: chest pain type other than typical angina to typical angina (through this part) \newline\textit{shape: categorical; highest to lowest: chest pain type other than typical angina, typical angina} & 18\% to 2\% & log-odds -2.04, OR 0.13; range across patients -2.26 to -1.57 & 220, 22 & train; test unclear (AUC gain 0.001 [-0.005, 0.008], n=61) & \textcolor{OliveGreen}{yes} & on the training data only \\
C32 & interaction: effect of number of vessels coloured by fluoroscopy (0 vessels to 2 vessels) by chest pain type: typical angina minus chest pain type other than typical angina; stronger in typical angina \newline\textit{shape: steady (equal steps); 100\% of patients same; verified; risk saturates near 0\%} & typical angina: $<$1\% to 55\%; chest pain type other than typical angina: 18\% to 99\% & log-odds difference 0.70, OR 2.00 & 22, 220 & train; test agrees (AUC gain 0.001 [-0.005, 0.008], n=61); group difference on test 0.70 [0.70, 0.70] & \textcolor{OliveGreen}{yes} & clearly \\
C33 & structure: its formula reproduces it on 242 training patients (largest difference 0.0) &  & log-odds 0.00, OR n/a & 242 & train & \textcolor{OliveGreen}{yes} &  \\
C34 & direction: ST depression on exercise: 0.0 mm to 2.8 mm (through this part) \newline\textit{shape: hinge at 1.0 mm; 100\% of patients same; fallback; one point with C37 (median fold 1.0 mm, typical patient 1.0 mm, 10th-90th percentile 1.0 mm to 1.0 mm)} & 18\% to 55\% & log-odds 0.87, OR 2.38; range across patients 0.00 to 1.74 & 242 & train; test unclear (AUC gain 0.005 [-0.005, 0.022], n=61) & \textcolor{OliveGreen}{yes} & on the training data only \\
C35 & direction: resting ECG: a normal resting ECG to resting ECG other than a normal resting ECG (through this part) \newline\textit{shape: categorical; highest to lowest: a normal resting ECG, resting ECG other than a normal resting ECG} & 99\% to 98\% (typical patient this part acts on) & log-odds -0.46, OR 0.63; range across patients -1.74 to 0.00 & 121, 121 & train; test disagrees (AUC gain -0.008 [-0.021, 0.000], n=61) & \textcolor{OliveGreen}{yes} & slightly; on the training data only \\
C36 & interaction: effect of ST depression on exercise (0.0 mm to 2.8 mm) by resting ECG: resting ECG other than a normal resting ECG minus a normal resting ECG; stronger in a normal resting ECG \newline\textit{shape: hinge at 1.0 mm; 100\% of patients same; fallback} & resting ECG other than a normal resting ECG: 18\% to 18\%; a normal resting ECG: 18\% to 55\% & log-odds difference -1.74, OR 0.18 & 121, 121 & train; test agrees (AUC gain -0.008 [-0.021, 0.000], n=61); group difference on test -1.74 [-1.74, -1.74] & \textcolor{OliveGreen}{yes} & clearly \\
\bottomrule
\end{tabularx}}
\par{\tiny\noindent
\begin{tabularx}{\linewidth}{@{}l>{\raggedright\arraybackslash}X>{\raggedright\arraybackslash}p{0.11\linewidth}>{\raggedright\arraybackslash}p{0.16\linewidth}>{\raggedright\arraybackslash}p{0.05\linewidth}>{\raggedright\arraybackslash}p{0.15\linewidth}l>{\raggedright\arraybackslash}p{0.08\linewidth}@{}}
\toprule
Claim & What is checked & Risk, typical patient & Estimate (log-odds; OR); range across patients & n & Test check & Verified & Words \\
\midrule
C37 & threshold: ST depression on exercise: fold at about 1.0 mm (starts) for patients with a normal resting ECG \newline\textit{shape: hinge at 1.0 mm; 100\% of patients same; fallback; one point with C37 (median fold 1.0 mm, typical patient 1.0 mm, 10th-90th percentile 1.0 mm to 1.0 mm)} & 13\% to 63\% (typical patient with a fold (patients with a normal resting ECG)) & fold 1.01 mm (IQR 1.01 to 1.01); past it log-odds 2.41, OR 11.1; range across patients 2.41 to 2.41 & 121, 41 & train; test unclear (AUC gain 0.005 [-0.005, 0.022], n=61) & \textcolor{OliveGreen}{yes} & on the training data only \\
C38 & subgroup: who this part acts on: 42 of 242 patients & 98\% to 99\% (typical patient this part acts on (risk without the part, then with it)) & share 0.17; log-odds 1.08, OR 2.95; range across patients 0.19 to 2.50 & 42, 200 & train; test unclear (AUC gain 0.005 [-0.006, 0.025], n=61) & \textcolor{OliveGreen}{yes} & slightly; on the training data only \\
C39 & structure: its formula reproduces it on 242 training patients (largest difference 2.22045e-16) &  & log-odds 2e-16, OR n/a & 242 & train & \textcolor{OliveGreen}{yes} &  \\
C40 & direction: number of vessels coloured by fluoroscopy: 0 vessels to 2 vessels (through this part) \newline\textit{shape: steady (equal steps); 100\% of patients same; verified} & 18\% to 36\% & log-odds 0.94, OR 2.56; range across patients 0.94 to 0.94 & 242 & train; test disagrees (AUC gain -0.002 [-0.019, 0.011], n=61) & \textcolor{OliveGreen}{yes} & on the training data only \\
C41 & direction: slope of the ST segment at peak exercise: slope of the ST segment at peak exercise other than a flat ST segment to a flat ST segment (through this part) \newline\textit{shape: categorical; highest to lowest: a flat ST segment, slope of the ST segment at peak exercise other than a flat ST segment} & 18\% to 23\% & log-odds 0.34, OR 1.41; range across patients 0.34 to 0.34 & 132, 110 & train; test disagrees (AUC gain -0.003 [-0.016, 0.005], n=61) & \textcolor{OliveGreen}{yes} & on the training data only \\
C42 & structure: its formula reproduces it on 242 training patients (largest difference 2.22045e-16) &  & log-odds 2e-16, OR n/a & 242 & train & \textcolor{OliveGreen}{yes} &  \\
C43 & direction: chest pain type: chest pain type other than non-anginal chest pain to non-anginal chest pain (through this part) \newline\textit{shape: categorical; highest to lowest: chest pain type other than non-anginal chest pain, non-anginal chest pain} & 18\% to 8\% & log-odds -0.92, OR 0.40; range across patients -0.92 to -0.92 & 172, 70 & train; test unclear (AUC gain 0.001 [-0.008, 0.013], n=61) & \textcolor{OliveGreen}{yes} & on the training data only \\
C44 & structure: its formula reproduces it on 242 training patients (largest difference 0.0) &  & log-odds 0.00, OR n/a & 242 & train & \textcolor{OliveGreen}{yes} &  \\
C45 & direction: maximum heart rate: 118 beats per minute to 175 beats per minute (through this part) \newline\textit{shape: hinge at 145 beats per minute; 100\% of patients same; verified; one point with C46 (median fold 145 beats per minute, typical patient 145 beats per minute, 10th-90th percentile 145 beats per minute to 145 beats per minute)} & 19\% to 14\% & log-odds -0.36, OR 0.70; range across patients -0.36 to -0.36 & 242 & train; test unclear (AUC gain 0.002 [0.000, 0.009], n=61) & \textcolor{OliveGreen}{yes} & slightly; on the training data only \\
C46 & threshold: maximum heart rate: fold at about 145 beats per minute (starts) \newline\textit{shape: hinge at 145 beats per minute; 100\% of patients same; verified; one point with C46 (median fold 145 beats per minute, typical patient 145 beats per minute, 10th-90th percentile 145 beats per minute to 145 beats per minute)} & 19\% to 13\% (typical patient with a fold) & fold 145 beats per minute (IQR 145 to 145); past it log-odds -0.44, OR 0.64; range across patients -0.44 to -0.44 & 242, 155 & train; test unclear (AUC gain 0.002 [0.000, 0.009], n=61) & \textcolor{OliveGreen}{yes} & on the training data only \\
C47 & structure: its formula reproduces it on 242 training patients (largest difference 0.0) &  & log-odds 0.00, OR n/a & 242 & train & \textcolor{OliveGreen}{yes} &  \\
\bottomrule
\end{tabularx}}
\paragraph{Not supported by the model.}
\noindent{\small Every claim passed its checks.}

\subsection{Coverage}\label{app:heart-s1-c1:coverage}
Coverage is computed by the current code (the function behind the \texttt{get\_coverage} tool) under these conventions. \emph{Candidate nodes} are all nodes of the explained model state except the output nodes. A candidate node is \emph{covered} when some annotation of the hierarchy (a leaf, or a composition together with its descendants) contains it and all its enabled outgoing edges. \emph{Exact Null} parts, which have no effect on any output for any input, count as covered: nodes with no path to an output over enabled, non-zero-weight edges (disconnected inputs and dead ends), and edges that have zero weight or touch such a node. Compositional coverage is the fraction of required compositions present: every annotation except a single global root containing every output must have a parent; a parentless, childless annotation made only of Null nodes is exempt.

\begin{itemize}
\item \textbf{Null, disconnected input (9):} \texttt{-10} (fbs:0), \texttt{-13} (restecg:1), \texttt{-14} (restecg:2), \texttt{-16} (exang:0), \texttt{-17} (exang:1), \texttt{-19} (slope:1), \texttt{-24} (thal:4), \texttt{-26} (thal:7), \texttt{-3} (sex:1).
\item \textbf{Uncovered nodes:} none.
\item \textbf{Uncovered edges:} none.
\item \textbf{Missing compositions:} none.
\end{itemize}

\subsection{Annotation hierarchy}\label{app:heart-s1-c1:hierarchy}
The hierarchy has 1 root; the global root is heart\_\allowbreak{}disease\_\allowbreak{}logit. Tags (A1, A2, \ldots) are used in the figures.
\begingroup\parindent=0pt\parskip=1pt
\hspace*{0.0em}$\triangleright$~A1: heart\_\allowbreak{}disease\_\allowbreak{}logit (composition, \cref{app:heart-s1-c1:ann-heart-disease-logit})\par
\hspace*{1.5em}$\triangleright$~A2: ca\_\allowbreak{}slope2\_\allowbreak{}relay (leaf, \cref{app:heart-s1-c1:ann-ca-slope2-relay})\par
\hspace*{1.5em}$\triangleright$~A3: cp3\_\allowbreak{}relay (leaf, \cref{app:heart-s1-c1:ann-cp3-relay})\par
\hspace*{1.5em}$\triangleright$~A4: st\_\allowbreak{}depression\_\allowbreak{}normal\_\allowbreak{}ecg\_\allowbreak{}gate (leaf, \cref{app:heart-s1-c1:ann-st-depression-normal-ecg-gate})\par
\hspace*{1.5em}$\triangleright$~A5: max\_\allowbreak{}heart\_\allowbreak{}rate\_\allowbreak{}hinge (leaf, \cref{app:heart-s1-c1:ann-max-heart-rate-hinge})\par
\hspace*{1.5em}$\triangleright$~A6: typical\_\allowbreak{}angina\_\allowbreak{}ca\_\allowbreak{}switch (leaf, \cref{app:heart-s1-c1:ann-typical-angina-ca-switch})\par
\endgroup

\subsection{Network figures}\label{app:heart-s1-c1:figures}
\begin{figure}[htbp]
\centering
\input{\explaneatcertdir{heart-s1-c1}heart-s1-c1-fig-full.tex}
\caption{heart-s1-c1: full view of the explained model state.}
\label{app:heart-s1-c1:fig-full}
\end{figure}
\begin{figure}[htbp]
\centering
\input{\explaneatcertdir{heart-s1-c1}heart-s1-c1-fig-level-0.tex}
\caption{heart-s1-c1: collapsed view at hierarchy level 0 (fully recomposed).}
\label{app:heart-s1-c1:fig-level-0}
\end{figure}
\begin{figure}[htbp]
\centering
\input{\explaneatcertdir{heart-s1-c1}heart-s1-c1-fig-level-1.tex}
\caption{heart-s1-c1: collapsed view at hierarchy level 1.}
\label{app:heart-s1-c1:fig-level-1}
\end{figure}

\subsection{Refactoring operations}\label{app:heart-s1-c1:operations}
The explained model state is the phenotype after these operations, replayed in order. All are function-preserving; the largest change in the model output over the 242 training rows is $3.33\times 10^{-16}$ (tolerance $1.00\times 10^{-9}$).
\begin{itemize}
\item \#26: split \texttt{-22} into \texttt{-22\_\allowbreak{}a}, \texttt{-22\_\allowbreak{}b}, \texttt{-22\_\allowbreak{}c}.
\item \#27: split \texttt{-20} into \texttt{-20\_\allowbreak{}a}, \texttt{-20\_\allowbreak{}b}.
\item \#28: split \texttt{-15} into \texttt{-15\_\allowbreak{}a}, \texttt{-15\_\allowbreak{}b}.
\end{itemize}

\subsection{Annotations}\label{app:heart-s1-c1:annotations}
\subsubsection{\texorpdfstring{A1: heart\_\allowbreak{}disease\_\allowbreak{}logit}{A1: heart disease logit}}\label{app:heart-s1-c1:ann-heart-disease-logit}
% annotation id ann_34, operation seq 34
\textbf{Kind:} Composition annotation (depth 0); children A2 (\cref{app:heart-s1-c1:ann-ca-slope2-relay}), A3 (\cref{app:heart-s1-c1:ann-cp3-relay}), A4 (\cref{app:heart-s1-c1:ann-st-depression-normal-ecg-gate}), A5 (\cref{app:heart-s1-c1:ann-max-heart-rate-hinge}), A6 (\cref{app:heart-s1-c1:ann-typical-angina-ca-switch}).

\begin{itemize}
\item \textbf{Entry nodes:} \texttt{-1} (age), \texttt{-2} (sex:0), \texttt{-3} (sex:1), \texttt{-5} (cp:2), \texttt{-7} (cp:4), \texttt{-8} (trestbps), \texttt{-9} (chol), \texttt{-10} (fbs:0), \texttt{-11} (fbs:1), \texttt{-13} (restecg:1), \texttt{-14} (restecg:2), \texttt{-15\_\allowbreak{}a} (thalach\_\allowbreak{}a), \texttt{-16} (exang:0), \texttt{-17} (exang:1), \texttt{-19} (slope:1), \texttt{-20\_\allowbreak{}a} (slope:2\_\allowbreak{}a), \texttt{-21} (slope:3), \texttt{-22\_\allowbreak{}a} (ca\_\allowbreak{}a), \texttt{-23} (thal:3), \texttt{-24} (thal:4), \texttt{-25} (thal:6), \texttt{-26} (thal:7).
\item \textbf{Exit nodes:} \texttt{0}.
\item \textbf{Internal nodes:} none.
\item \textbf{Edges (weight):} \texttt{-1}$\to$\texttt{0}~(0.399), \texttt{-2}$\to$\texttt{0}~(-1.87), \texttt{-5}$\to$\texttt{0}~(0.443), \texttt{-7}$\to$\texttt{0}~(1.25), \texttt{-8}$\to$\texttt{0}~(0.665), \texttt{-9}$\to$\texttt{0}~(0.931), \texttt{-11}$\to$\texttt{0}~(-0.315), \texttt{-21}$\to$\texttt{0}~(0.197), \texttt{-23}$\to$\texttt{0}~(-0.866), \texttt{-25}$\to$\texttt{0}~(-0.861), \texttt{664}$\to$\texttt{0}~(0.92), \texttt{1173}$\to$\texttt{0}~(-1.24), \texttt{1801}$\to$\texttt{0}~(1.4), \texttt{2009}$\to$\texttt{0}~(-0.151), \texttt{2245}$\to$\texttt{0}~(1.53), \texttt{2650}$\to$\texttt{0}~(-0.172), \texttt{3210}$\to$\texttt{0}~(-0.833), \texttt{-22\_\allowbreak{}a}$\to$\texttt{0}~(1.3), \texttt{-20\_\allowbreak{}a}$\to$\texttt{0}~(1.72), \texttt{-15\_\allowbreak{}a}$\to$\texttt{0}~(0.00188).
\item \textbf{With its descendants:} 38 nodes.
\item \textbf{Covers} (coverage, own and as part of the hierarchy): 37 candidate nodes.
\end{itemize}

\textbf{Hypothesis.} Hypothesis: the heart-disease logit is essentially a logistic regression. 13 direct input-$>$output edges (largest: sex:0 -1.87, slope:2 +1.72, ca +1.30, cp:4 +1.25, chol +0.93, thal:3 -0.87, thal:6 -0.86, trestbps +0.67, cp:2 +0.44, age +0.40, fbs:1 -0.32, slope:3 +0.20, thalach +0.002) are added to five children: ca\_\allowbreak{}slope2\_\allowbreak{}relay and the 2245 relay (raise the effective ca weight to 2.60 per SD), cp3\_\allowbreak{}relay (non-anginal pain lowers risk), st\_\allowbreak{}depression\_\allowbreak{}normal\_\allowbreak{}ecg\_\allowbreak{}gate (ST depression raises risk only with a normal resting ECG), typical\_\allowbreak{}angina\_\allowbreak{}ca\_\allowbreak{}switch's 3210 (typical angina lowers risk, less with more vessels) and max\_\allowbreak{}heart\_\allowbreak{}rate\_\allowbreak{}hinge (small). Nine inputs are unconnected (sex:1, fbs:0, restecg:1, restecg:2, exang:0, exang:1, slope:1, thal:4, thal:7): exercise-induced angina is ignored. Evidence (train, n=242): AUC 0.9392; all hidden units replaced by their best linear fits 0.9355 (= L2 logistic regression 0.9355); all ReLUs as identity 0.9078; all hidden mean-ablated 0.9241. The gain over linear (+0.0037) comes from the ECG/ST gate (+0.0034 restored alone), the heart-rate hinge (+0.0008) and the typical-angina switch (+0.0003); the relays contribute none. Direct edges carry logit SD 4.25 of 5.63.

\textbf{Closed form} (children as named functions):
{\small\begingroup\allowdisplaybreaks
\begin{align*}
y_{0} &= \sigma\bigl(z_{0}\bigr) \\
z_{0} &= -0.276 + 1.3 x_{ca\_a} + 0.666 x_{trestbps} \\
&\qquad + 0.92 \operatorname{ca\_slope2\_relay}\bigl(x_{ca\_b}, x_{slope:2\_b}\bigr) \\
&\qquad + 0.00188 x_{thalach\_a} + 1.25 x_{cp:4} \\
&\qquad + 1.4 \operatorname{st\_depression\_normal\_ecg\_gate}\bigl(x_{restecg:0}, x_{oldpeak}\bigr) \\
&\qquad + 1.53 \operatorname{typical\_angina\_ca\_switch\_0}\bigl(x_{ca\_c}, x_{cp:1}\bigr) \\
&\qquad + 0.197 x_{slope:3} + 0.399 x_{age} + 1.72 x_{slope:2\_a} + 0.443 x_{cp:2} + 0.931 x_{chol} \\
&\qquad - 1.87 x_{sex:0} - 1.24 \operatorname{cp3\_relay}\bigl(x_{cp:3}\bigr) - 0.315 x_{fbs:1} \\
&\qquad - 0.861 x_{thal:6} - 0.151 \operatorname{max\_heart\_rate\_hinge\_0}\bigl(x_{thalach\_b}\bigr) \\
&\qquad - 0.172 \operatorname{max\_heart\_rate\_hinge\_1}\bigl(x_{thalach\_b}\bigr) \\
&\qquad - 0.833 \operatorname{typical\_angina\_ca\_switch\_1}\bigl(x_{ca\_c}, x_{cp:1}\bigr) \\
&\qquad - 0.866 x_{thal:3}
\end{align*}\endgroup}

Expanded over the children:
{\small\begingroup\allowdisplaybreaks
\begin{align*}
y_{0} &= \sigma\bigl(z_{0}\bigr) \\
z_{0} &= -0.276 + 1.3 x_{ca\_a} + 0.666 x_{trestbps} + 0.92 \operatorname{relu}\bigl(0.783 \\
&\qquad + 0.185 x_{slope:2\_b} + 0.456 x_{ca\_b}\bigr) + 0.00188 x_{thalach\_a} + 1.25 x_{cp:4} \\
&\qquad + 1.4 \operatorname{relu}\bigl(-1.52 + 1.54 x_{restecg:0} + 0.818 x_{oldpeak}\bigr) \\
&\qquad + 1.53 \operatorname{relu}\bigl(1.28 + 0.578 x_{ca\_c}\bigr) + 0.197 x_{slope:3} \\
&\qquad + 0.399 x_{age} + 1.72 x_{slope:2\_a} + 0.443 x_{cp:2} + 0.931 x_{chol} - 1.87 x_{sex:0} \\
&\qquad - 1.24 \operatorname{relu}\bigl(1.82 + 0.336 x_{cp:3}\bigr) - 0.315 x_{fbs:1} \\
&\qquad - 0.861 x_{thal:6} - 0.151 \operatorname{relu}\bigl(0.201 + 0.833 x_{thalach\_b}\bigr) \\
&\qquad - 0.172 \operatorname{relu}\bigl(0.859 + 0.982 \operatorname{relu}\bigl(0.201 \\
&\qquad + 0.833 x_{thalach\_b}\bigr)\bigr) - 0.833 \operatorname{relu}\bigl(0.42 + 0.901 x_{cp:1} \\
&\qquad - 0.644 \operatorname{relu}\bigl(1.28 + 0.578 x_{ca\_c}\bigr)\bigr) - 0.866 x_{thal:3}
\end{align*}\endgroup}

where $y_{0}$ is node \texttt{0}; $z_{i}$, where shown, is the pre-activation of a sigmoid output $y_{i}$; weights and biases are shown to 3 significant figures.
 On the training split (n = 242), the displayed closed form differs from the forward pass by at most 0.00427, and the exact function by at most $4.44\times 10^{-16}$.

\textbf{Evidence} (2 items; full payloads in the JSON record):
\begin{itemize}
\item \textsc{analytical}, \texttt{analytical\_\allowbreak{}formula}; data split: n/a (analytical). Root closed form (get\_\allowbreak{}formula force=True) and effective linear weights through the always-on relays
\item \textsc{empirical}, \texttt{ablation\_\allowbreak{}study}; data split: train. Train split, full sample; ablations and linear replacements on in-memory copies (no graph edits)
\end{itemize}

\textbf{Authorship:} drafted by AI-assisted draft on 2026-11-10; verified by nobody (unverified).

\subsubsection{\texorpdfstring{A2: ca\_\allowbreak{}slope2\_\allowbreak{}relay}{A2: ca slope2 relay}}\label{app:heart-s1-c1:ann-ca-slope2-relay}
% annotation id ann_29, operation seq 29
\textbf{Kind:} Leaf annotation (depth 1); parent A1 (\cref{app:heart-s1-c1:ann-heart-disease-logit}).

\begin{itemize}
\item \textbf{Entry nodes:} \texttt{-22\_\allowbreak{}b} (ca\_\allowbreak{}b), \texttt{-20\_\allowbreak{}b} (slope:2\_\allowbreak{}b).
\item \textbf{Exit nodes:} \texttt{664}.
\item \textbf{Internal nodes:} none.
\item \textbf{Edges (weight):} \texttt{-22\_\allowbreak{}b}$\to$\texttt{664}~(0.456), \texttt{-20\_\allowbreak{}b}$\to$\texttt{664}~(0.185).
\item \textbf{Covers} (coverage, own and as part of the hierarchy): 2 candidate nodes.
\end{itemize}

\textbf{Hypothesis.} Hypothesis: node 664 = ReLU(0.456 ca + 0.185 slope:2 + 0.783) is an affine relay that strengthens the direct ca and flat-ST-slope (slope:2) effects. Evidence (train, n=242): pre-activation 0.28..2.19, never clips, so the part is exactly affine on the data; it adds +0.42 logit per SD of ca and +0.17 per SD of slope:2 (output weight 0.92). Mean ablation AUC 0.9392 -$>$ 0.9375; linear replacement leaves AUC at 0.9392. No concept beyond reweighting.

\textbf{Closed form:}
{\small\begingroup\allowdisplaybreaks
\begin{align*}
y_{0} &= \operatorname{relu}\bigl(0.783 + 0.185 x_{slope:2\_b} + 0.456 x_{ca\_b}\bigr)
\end{align*}\endgroup}

where $y_{0}$ is node \texttt{664}; $z_{i}$, where shown, is the pre-activation of a sigmoid output $y_{i}$; weights and biases are shown to 3 significant figures.
 On the training split (n = 242), the displayed closed form differs from the forward pass by at most $4.47\times 10^{-4}$, and the exact function by at most 0.

\textbf{Evidence} (2 items; full payloads in the JSON record):
\begin{itemize}
\item \textsc{analytical}, \texttt{analytical\_\allowbreak{}formula}; data split: train (inferred, inferred from payload keys). Closed form from get\_\allowbreak{}formula (force=True)
\item \textsc{empirical}, \texttt{empirical\_\allowbreak{}viz} (heatmap); data split: train. Train split, full sample (242); local forward passes on in-memory copies for ablations
\end{itemize}

\textbf{Authorship:} drafted by AI-assisted draft on 2026-11-10; verified by nobody (unverified).

\subsubsection{\texorpdfstring{A3: cp3\_\allowbreak{}relay}{A3: cp3 relay}}\label{app:heart-s1-c1:ann-cp3-relay}
% annotation id ann_30, operation seq 30
\textbf{Kind:} Leaf annotation (depth 1); parent A1 (\cref{app:heart-s1-c1:ann-heart-disease-logit}).

\begin{itemize}
\item \textbf{Entry nodes:} \texttt{-6} (cp:3).
\item \textbf{Exit nodes:} \texttt{1173}.
\item \textbf{Internal nodes:} none.
\item \textbf{Edges (weight):} \texttt{-6}$\to$\texttt{1173}~(0.336).
\item \textbf{Covers} (coverage, own and as part of the hierarchy): 1 candidate node.
\end{itemize}

\textbf{Hypothesis.} Hypothesis: node 1173 = ReLU(0.336 cp:3 + 1.818) is an always-on relay through which non-anginal chest pain (cp:3) lowers predicted risk. Evidence (train, n=242): pre-activation 1.60..2.34, never clips (entry-exit correlation 1.0); with output weight -1.24 it contributes -0.417 logit per SD of cp:3, about -0.92 logit for non-anginal pain vs other types. cp:3 has no direct edge, so this is its only route. Mean ablation AUC 0.9373; linear replacement 0.9392.

\textbf{Closed form:}
{\small\begingroup\allowdisplaybreaks
\begin{align*}
y_{0} &= \operatorname{relu}\bigl(1.82 + 0.336 x_{cp:3}\bigr)
\end{align*}\endgroup}

where $y_{0}$ is node \texttt{1173}; $z_{i}$, where shown, is the pre-activation of a sigmoid output $y_{i}$; weights and biases are shown to 3 significant figures.
 On the training split (n = 242), the displayed closed form differs from the forward pass by at most 0.00246, and the exact function by at most 0.

\textbf{Evidence} (2 items; full payloads in the JSON record):
\begin{itemize}
\item \textsc{analytical}, \texttt{analytical\_\allowbreak{}formula}; data split: train (inferred, inferred from payload keys). Closed form from get\_\allowbreak{}formula (force=True)
\item \textsc{empirical}, \texttt{empirical\_\allowbreak{}viz} (line\_\allowbreak{}plot); data split: train. Train split, full sample (242); local forward passes on in-memory copies for ablations
\end{itemize}

\textbf{Authorship:} drafted by AI-assisted draft on 2026-11-10; verified by nobody (unverified).

\subsubsection{\texorpdfstring{A4: st\_\allowbreak{}depression\_\allowbreak{}normal\_\allowbreak{}ecg\_\allowbreak{}gate}{A4: st depression normal ecg gate}}\label{app:heart-s1-c1:ann-st-depression-normal-ecg-gate}
% annotation id ann_31, operation seq 31
\textbf{Kind:} Leaf annotation (depth 1); parent A1 (\cref{app:heart-s1-c1:ann-heart-disease-logit}).

\begin{itemize}
\item \textbf{Entry nodes:} \texttt{-12} (restecg:0), \texttt{-18} (oldpeak).
\item \textbf{Exit nodes:} \texttt{1801}.
\item \textbf{Internal nodes:} none.
\item \textbf{Edges (weight):} \texttt{-12}$\to$\texttt{1801}~(1.54), \texttt{-18}$\to$\texttt{1801}~(0.817).
\item \textbf{Covers} (coverage, own and as part of the hierarchy): 2 candidate nodes.
\end{itemize}

\textbf{Hypothesis.} Hypothesis: node 1801 = ReLU(1.542 restecg:0 + 0.817 oldpeak - 1.521) is a gate in which ST depression (oldpeak) raises predicted risk only for patients with a normal resting ECG (restecg:0), once oldpeak is above about its mean (threshold -0.03 SD); with an abnormal resting ECG the threshold is 3.75 SD and the gate is effectively closed. Evidence (train, n=242): active for 41 of 121 normal-ECG patients and 1 of 121 others; adds 1.40 x activation = +1.14 logit per SD oldpeak when open. It is the model's main nonlinearity: mean ablation AUC 0.9339, linear replacement 0.9362, and restoring it onto the all-linear model gains +0.0034 AUC. oldpeak has no direct edge, so the model ignores ST depression for patients with an abnormal resting ECG.

\textbf{Closed form:}
{\small\begingroup\allowdisplaybreaks
\begin{align*}
y_{0} &= \operatorname{relu}\bigl(-1.52 + 1.54 x_{restecg:0} + 0.818 x_{oldpeak}\bigr)
\end{align*}\endgroup}

where $y_{0}$ is node \texttt{1801}; $z_{i}$, where shown, is the pre-activation of a sigmoid output $y_{i}$; weights and biases are shown to 3 significant figures.
 On the training split (n = 242), the displayed closed form differs from the forward pass by at most 0.00514, and the exact function by at most 0.

\textbf{Evidence} (2 items; full payloads in the JSON record):
\begin{itemize}
\item \textsc{analytical}, \texttt{analytical\_\allowbreak{}formula}; data split: train (inferred, inferred from payload keys). Closed form from get\_\allowbreak{}formula (force=True)
\item \textsc{empirical}, \texttt{empirical\_\allowbreak{}viz} (heatmap); data split: train. Train split, full sample (242); local forward passes on in-memory copies for ablations
\end{itemize}

\textbf{Authorship:} drafted by AI-assisted draft on 2026-11-10; verified by nobody (unverified).

\subsubsection{\texorpdfstring{A5: max\_\allowbreak{}heart\_\allowbreak{}rate\_\allowbreak{}hinge}{A5: max heart rate hinge}}\label{app:heart-s1-c1:ann-max-heart-rate-hinge}
% annotation id ann_32, operation seq 32
\textbf{Kind:} Leaf annotation (depth 1); parent A1 (\cref{app:heart-s1-c1:ann-heart-disease-logit}).

\begin{itemize}
\item \textbf{Entry nodes:} \texttt{-15\_\allowbreak{}b} (thalach\_\allowbreak{}b).
\item \textbf{Exit nodes:} \texttt{2009}, \texttt{2650}.
\item \textbf{Internal nodes:} none.
\item \textbf{Edges (weight):} \texttt{-15\_\allowbreak{}b}$\to$\texttt{2009}~(0.833), \texttt{2009}$\to$\texttt{2650}~(0.982).
\item \textbf{Covers} (coverage, own and as part of the hierarchy): 1 candidate node.
\end{itemize}

\textbf{Hypothesis.} Hypothesis: a small hinge on maximum heart rate (thalach). Node 2009 = ReLU(0.833 thalach + 0.201) feeds the output (-0.151) and the always-on relay 2650 = ReLU(0.982*2009 + 0.859) (-0.172), a net -0.32 x 2009: higher maximum heart rate lowers predicted risk, but only above about -0.24 SD; below that it has no effect (the direct thalach edge, +0.002, is negligible). Evidence (train, n=242): 2009 active for 64\% of patients; logit contribution SD 0.16 (small); mean ablation AUC 0.9391, linear replacement 0.9392, restore-onto-linear gain +0.0008.

\textbf{Closed form:}
{\small\begingroup\allowdisplaybreaks
\begin{align*}
y_{0} &= \operatorname{relu}\bigl(0.201 + 0.833 x_{thalach\_b}\bigr) \\
y_{1} &= \operatorname{relu}\bigl(0.859 + 0.982 \operatorname{relu}\bigl(0.201 \\
&\qquad + 0.833 x_{thalach\_b}\bigr)\bigr)
\end{align*}\endgroup}

where $y_{0}$ is node \texttt{2009}; $y_{1}$ is node \texttt{2650}; $z_{i}$, where shown, is the pre-activation of a sigmoid output $y_{i}$; weights and biases are shown to 3 significant figures.
 On the training split (n = 242), the displayed closed form differs from the forward pass by at most $7.41\times 10^{-4}$, and the exact function by at most 0.

\textbf{Evidence} (2 items; full payloads in the JSON record):
\begin{itemize}
\item \textsc{analytical}, \texttt{analytical\_\allowbreak{}formula}; data split: train (inferred, inferred from payload keys). Closed form from get\_\allowbreak{}formula (force=True)
\item \textsc{empirical}, \texttt{empirical\_\allowbreak{}viz} (line\_\allowbreak{}plot); data split: train. Train split, full sample (242); local forward passes on in-memory copies for ablations
\end{itemize}

\textbf{Authorship:} drafted by AI-assisted draft on 2026-11-10; verified by nobody (unverified).

\subsubsection{\texorpdfstring{A6: typical\_\allowbreak{}angina\_\allowbreak{}ca\_\allowbreak{}switch}{A6: typical angina ca switch}}\label{app:heart-s1-c1:ann-typical-angina-ca-switch}
% annotation id ann_33, operation seq 33
\textbf{Kind:} Leaf annotation (depth 1); parent A1 (\cref{app:heart-s1-c1:ann-heart-disease-logit}).

\begin{itemize}
\item \textbf{Entry nodes:} \texttt{-22\_\allowbreak{}c} (ca\_\allowbreak{}c), \texttt{-4} (cp:1).
\item \textbf{Exit nodes:} \texttt{2245}, \texttt{3210}.
\item \textbf{Internal nodes:} none.
\item \textbf{Edges (weight):} \texttt{-22\_\allowbreak{}c}$\to$\texttt{2245}~(0.578), \texttt{-4}$\to$\texttt{3210}~(0.901), \texttt{2245}$\to$\texttt{3210}~(-0.644).
\item \textbf{Covers} (coverage, own and as part of the hierarchy): 2 candidate nodes.
\end{itemize}

\textbf{Hypothesis.} Hypothesis: (i) node 2245 = ReLU(0.578 ca + 1.276) is an always-on ca relay adding +0.88 logit per SD of ca (output weight 1.528); (ii) node 3210 = ReLU(0.901 cp:1 - 0.644*2245 + 0.420) is a typical-angina switch: it is on only for cp:1 patients and then lowers the logit by 0.833 x its activation, from about -2.26 at the lowest ca (-0.73 SD) to about -1.22 at the highest ca (2.64 SD). So typical angina lowers predicted risk, and the reduction shrinks with more narrowed vessels. Evidence (train, n=242): 2245 never clips (0.86..2.80); 3210 active for all 22 cp:1 patients and none of the 220 others. Mean ablation AUC 2245 0.9339, 3210 0.9354; linear replacement 0.9392 / 0.9391 (3210 is a cp:1 switch whose within-group response is linear in ca).

\textbf{Closed form:}
{\small\begingroup\allowdisplaybreaks
\begin{align*}
y_{0} &= \operatorname{relu}\bigl(1.28 + 0.578 x_{ca\_c}\bigr) \\
y_{1} &= \operatorname{relu}\bigl(0.42 + 0.901 x_{cp:1} - 0.644 \operatorname{relu}\bigl(1.28 \\
&\qquad + 0.578 x_{ca\_c}\bigr)\bigr)
\end{align*}\endgroup}

where $y_{0}$ is node \texttt{2245}; $y_{1}$ is node \texttt{3210}; $z_{i}$, where shown, is the pre-activation of a sigmoid output $y_{i}$; weights and biases are shown to 3 significant figures.
 On the training split (n = 242), the displayed closed form differs from the forward pass by at most 0.00423, and the exact function by at most 0.

\textbf{Evidence} (2 items; full payloads in the JSON record):
\begin{itemize}
\item \textsc{analytical}, \texttt{analytical\_\allowbreak{}formula}; data split: train (inferred, inferred from payload keys). Closed form from get\_\allowbreak{}formula (force=True)
\item \textsc{empirical}, \texttt{empirical\_\allowbreak{}viz} (heatmap); data split: train. Train split, full sample (242); local forward passes on in-memory copies for ablations
\end{itemize}

\textbf{Authorship:} drafted by AI-assisted draft on 2026-11-10; verified by nobody (unverified).

\subsection{Audit against the definitions}\label{app:heart-s1-c1:audit}
Automatic findings on the stored explanation (codes as in the JSON record's \texttt{issues}); they flag what an examiner should check, not verdicts.
\begin{itemize}
\item \textsc{warning} \texttt{unverified}: No human verification is recorded for this explanation.
\end{itemize}
\end{drafted}
